\section{Finding 1: sHash Cannot Be Searched Correctly}
\label{sec:shash}

\subsection{How sHash compares two images}

In the paper's storage design (its Fig.~11b) each sHash BK-tree node holds one image's list of
piece hashes. Comparing image $A$, with pieces $a_1,\dots,a_m$, to image $B$, with pieces
$b_1,\dots,b_n$, works as follows: for each piece of $A$, find its closest piece in $B$, then
average those closest distances:
\[
  d_s(A \to B) \;=\; \frac{1}{m}\sum_{i=1}^{m}\;\min_{j}\; d_H(a_i, b_j).
\]
The averaging is not stated in the paper; it comes from the authors' code and data, which our
C11 reimplementation reproduces exactly. Only the pieces of the first image have to find a
match, so the comparison runs one way.

\subsection{It is not symmetric}

Let image A have 3 pieces and image B the same 3 plus one more that matches nothing in A
(Figure~\ref{fig:asym}). With matching pieces 2 apart and the unmatched piece 34 from its
nearest neighbour:
\[
  d_s(A\to B) = \tfrac{2+2+2}{3} = 2, \qquad d_s(B\to A) = \tfrac{2+2+2+34}{4} = 10.
\]
The same two images get two different answers. The data the authors shared with us shows this
in practice. The paper's text describes comparing the \emph{new} image to the original, but
their reported distances match the original$\to$copy direction on all 1,802 pairs and the
copy$\to$original direction on only 240. With a one-way comparison, this is an easy detail to
miss.

\begin{figure}[ht]
\centering
\begin{tikzpicture}[font=\small, >={Stealth[length=2mm]},
  pc/.style={draw=ink2, fill=orbblue!20, minimum size=6mm, rounded corners=1pt},
  px/.style={draw=ink2, fill=shashorange!35, minimum size=6mm, rounded corners=1pt}]
  \node at (0,1.0) {\textbf{Image A} (3 pieces)};
  \foreach \i in {0,1,2} \node[pc] (a\i) at (\i*0.9-0.9,0.2) {};
  \node at (5.2,1.0) {\textbf{Image B} (4 pieces)};
  \foreach \i in {0,1,2} \node[pc] (b\i) at (\i*0.9+3.85,0.2) {};
  \node[px] (b3) at (6.55,0.2) {};
  \foreach \i in {0,1,2} \draw[<->, ink2] (a\i.south) to[bend right=25] node[pos=0.5, below, font=\footnotesize] {2} (b\i.south);
  \node[font=\footnotesize, text=shashorange!80!black] at (6.55,-0.55) {no match: 34};
  \node[align=left, anchor=west] at (8.0,0.45) {$A\to B = 2$};
  \node[align=left, anchor=west] at (8.0,-0.15) {$B\to A = 10$};
\end{tikzpicture}
\caption{The sHash comparison is one-way. B's extra piece costs nothing when A is compared to
B, and a lot when B is compared to A.}
\label{fig:asym}
\end{figure}

\subsection{The triangle inequality fails}

Let $A$ be an image of a cat, $C$ an image of a dog, and $B$ an image containing both. Take
cat-to-cat and dog-to-dog pieces as 2 apart and cat-to-dog as 30 (illustrative values):
\[
  d(A,B) = 2, \qquad d(B,C) = \tfrac{30+2}{2} = 16, \qquad d(A,C) = 30 \;>\; 2 + 16 = 18.
\]
The direct distance is longer than the detour through $B$, which a metric does not allow. A
composite image sits close to two unrelated images at once.

\subsection{Why this loses real copies}

The tree's pruning relies on exactly that inequality. Suppose the composite $B$ is stored at
a node, and a genuine copy $X$ of a cat is stored below it. Because of $B$'s dog, the stored
edge is $d(B,X)=16$. A cat query $q$ arrives with $d(q,B)=2$, and the true copy is only
$d(q,X)=2$ away. With threshold $t=5$ the search descends only into edges in
$[2-5,\,2+5]=[-3,7]$, so the edge labelled 16 is skipped (Figure~\ref{fig:prune}). The search
reports ``no copy'' and nothing looks wrong. A silently missed copy is the most costly error
for this system.

\begin{figure}[ht]
\centering
\begin{tikzpicture}[font=\small, >={Stealth[length=2mm]},
  n/.style={draw=ink2, circle, minimum size=9mm, align=center, fill=white}]
  \node[n] (B) at (0,0) {B\\[-3pt]\footnotesize cat+dog};
  \node[n, fill=rule!60] (X) at (0,-2.4) {X\\[-3pt]\footnotesize cat};
  \node[n, draw=orbblue, very thick] (q) at (-4.2,0) {q\\[-3pt]\footnotesize cat};
  \draw[->] (B) -- node[right] {stored edge 16} (X);
  \draw[->, dashed, ink2] (q) -- node[above] {$d(q,B)=2$} (B);
  \draw[->, dashed, orbblue] (q) -- node[below left] {true $d(q,X)=2$} (X);
  \node[anchor=west, align=left, font=\footnotesize] at (2.4,-0.6)
    {threshold $t=5$\\ search visits edges in $[-3,\,7]$\\ edge 16 is skipped:\\ \textbf{the real copy is never checked}};
\end{tikzpicture}
\caption{A BK-tree search missing a real copy because the sHash comparison breaks the triangle
inequality.}
\label{fig:prune}
\end{figure}

We shared this analysis with the authors, who replied that it may help explain why the tree's
results are not perfect. Our claim is limited to what the code and the mathematics show: the sHash
comparison is not a metric, so standard BK-tree pruning can miss a valid match. This changed
the project. It made little sense to accelerate a search structure that can miss real copies,
so we turned to replacing the signal.
